\documentclass{article}
\usepackage{graphicx}
\usepackage{epsfig}
\usepackage{amsmath}
\usepackage{amsfonts}
\usepackage{amssymb}
\usepackage{latexsym}
\usepackage{mathrsfs}
\usepackage{algpseudocode}

\oddsidemargin -0.0 in \topmargin -25 mm \textheight 10in
\textwidth 6.5 in

%\pagenumbering{empty}
\begin{document}
\begin{centering}

\textbf{\large Assignment 5 CSE 16-02: Applied Discrete Mathematics}\\


\textbf{\large Winter 2026}\\
Instructor: Cedric Westphal  \\
Due February 12th, 2026 11:59 PM\\

\end{centering}

\vspace{12pt}
\noindent Complete the questions carefully and show your work. Please remember to:
\begin{itemize}
	\item Scan the pages and combine them into a single PDF file. Make sure your work is readable and that pages are oriented correctly.
	\item Upload your PDF file to Canvas/Gradescope under Assignment 5.
	\item In Gradescope select the page(s) for each problem.
\end{itemize}

\noindent Notation: $\emptyset$ is the empty set, $\mathbb{N}$ is the set of natural numbers, $\mathbb{Z}$ is the set of integers, $\mathbb{Q}$ is the set of rational numbers, and $\mathbb{R}$ is the set of real numbers. $\mathscr{P}(A)$ is the power set of $A$.\\
\textbf{[46 points total in assignment.]}

\begin{enumerate}
\setcounter{enumi}{-1} 
\item Each answer is clearly written and each answer is correctly mapped to its rubric on Gradescope. [4 points]

\item  For $a,b \in \mathbb{N}$, prove the following statements below. Hint: to prove that $g$ is the gcd($a,b$) you need to prove that $g | a$ and $g|b$ and that $g$ is the greatest such common divisor.\textbf{[10 points]}
\begin{enumerate}
    \item Let $a = 5x$ and $b = 5y+7$ for some $x$ and $y$ in $\mathbb{N}$. Prove that gcd($a,b$) = gcd($x,b$). Write a proof that shows: (i) that gcd($a,b)| x$, then (ii) show that it implies that gcd($a,b$) equals gcd($x,b$). \textbf{[2 points]}
    \item If both $a$ and $b$ are divisible by 4, then gcd($a,b$) = 4 gcd($a/4,b/4$). \textbf{[2 points]}
    \item If $m$ is a common divisor of $a$ and $b$, show the general case of the previous question, namely that gcd($a,b$) = $m$ gcd($a/m,b/m$). \textbf{[2 points]}
    \item Prove that gcd($a,b$) = gcd($a+b,b$). \textbf{[2 points]}
    \item Let \(a\)  and \(b\)  be integers, not both zero. Prove that \(\gcd ( a , b )\)  is the smallest positive integer that can be written in the form \(ax+by\) for integers \(x\)  and \(y\) .
    % \item Is it true or false that, for $a>b$, 2gcd($a,b$) = gcd($a+b,a-b$)? Give a clear explanation of your reasoning. \textbf{[2 points]}
\end{enumerate}

\textbf{Sol}
\begin{enumerate}
    \item \begin{enumerate}
        \item Since \(5y+7\) is not divisible by 5, the gcd of \(a\) and \(b\) must divide \(a=5x\), by Euclid's Lemma, the gcd must divide \(x\). 
        \item \textbf{Forward Direction}: Suppose $d$ is a common divisor of $a$ and $b$ ($d \mid a$ and $d \mid b$). From part (i), we already proved that if $d$ divides both $5x$ and $b$, then $d \mid x$. Thus, $d$ is a common divisor of $x$ and $b$. \textbf{Reverse Direction}: Suppose $d$ is a common divisor of $x$ and $b$ ($d \mid x$ and $d \mid b$). If $d \mid x$, then $d$ must also divide any multiple of $x$, specifically $d \mid 5x$. Since $a = 5x$, $d$ divides $a$. Thus, $d$ is a common divisor of $a$ and $b$.
    \end{enumerate} 
    \item Since \(a\) and \(b \) are divisible by \(4\), there exists \(a^\prime,b^\prime\) such that \(a = 4a^\prime, b = 4b^\prime\). We want to show there exists some \(g^\prime\) such that \(g = \text{gcd}(a,b) = 4g^\prime\). Let \(d\) be the gcd of \(a^\prime\) and \(b^\prime\). We then know that \(4d\mid a\) and \(4d\mid b\). Since \(4\) is a common factor of \(a\) and \(b\), the gcd of \(a\) and \(b\) can be written as \(4c^\prime\). Chaining everything together we get 
    \begin{equation}
        \gcd(a,b) = 4c^\prime = 4d = 4\gcd(a^\prime,b^\prime) = 4\gcd \left(\frac{a}{4},\frac{b}{4}\right)
    \end{equation}
    \item The proof is identical but replace 4 with \(m\). 
    \item 
    (\(\leq\)) 
    
    Let \(x=\gcd(a,b)\). Since \(x\) divides \(a\) and \(b, \) it must also divide \(a+b\) (note \(a+b = x(a^\prime + b^\prime)\) for some \(a^\prime, b^\prime \) by definition of divides). This means that \(\gcd(a+b,b)\) must be at least \(\gcd(a,b)\). 

    (\(\geq\))

    Let \(y=\gcd(a+b,b)\). \(y\) must also divide \(a+b-b\) (let \(\alpha_1 y = a+b\) and \(\alpha_2 y = b\). This means \((\alpha_1-\alpha_2 )y = a + b - b\)). This means that \(\gcd(a,b)\) must be at least \(\gcd(a+b,b)\). 

    Since they are \(\geq\) and \(\leq\), they must be equal. 
    \item Consider the set of all positive numbers that \(ax+by\) form given that \(a,b\) are nonzero and \(x,y\) are in the integers. This set is nonempty since at least one of the two \(a,b\) is nonzero. By the Well Ordering Principle, there must exist a least element. Call this least element \(d\). By the division algorithm, we know that 
    \begin{equation}
        \begin{split}
            a &= dq_a+r_a\\
            b &= dq_b+r_b
        \end{split}
    \end{equation}
    for some \(q_a,q_b\in \mathbb Z\) and \(r_a,r_b\in\{0,\dots, d-1\}\). We will work with \(r=r_a, q=q_a\) and the proof proceeds without loss of generality with \(b\). 
    $$r = a - dq$$$$r = a - (ax + by)q$$ $$r = a(1 - xq) + b(-yq)$$ 
    This shows that $r$ is also a linear combination of $a$ and $b$. If $r > 0$, then $r$ would be an element of $S$. However, we stated that $d$ is the smallest positive element of $S$, and the division algorithm tells us $r < d$. The only way to avoid a contradiction is if $r = 0$. Therefore, $d \mid a$. By the same logic, $d \mid b$. Thus, $d$ is a common divisor of $a$ and $b$.
\end{enumerate}



\item We will prove that if $n$ is an odd integer $\in \mathbb{N}$, then 16 divides $n^4 - 1$. We're going to do this in two ways. \textbf{[7 points]} 
\begin{enumerate}
    \item First, prove that if $n$ is an odd integer $\in \mathbb{N}$, then 8 divides $n^2 - 1$.
    \item Prove that if $n$ is an odd integer, then $n^2+1$ is even.
    \item Using the identity $n^4 - 1 = (n^2 - 1)(n^2 + 1)$ and your results from steps 1 and 2, prove that $16$ divides $n^4 - 1$.
    
    \item Let $n = 2k + 1$. Express $n^4 - 1$ using the binomial theorem as a sum of terms of the form $a_i k^i$. Write this sum explicitly.
    \item Consider the last two terms of your sum, $a_2 k^2$ and $a_1 k$. Rewrite $a_2 k^2 + a_1 k$ in the form $b k^2 + c(k^2 + k)$ where $b$ is an integer.
    \item Show that $b$ and \(c(k^2+k)\) is divisible by 8.
    \item Prove the result stated at the beginning of this Problem, using the previous three steps. 
\end{enumerate}

\textbf{sol}
\begin{enumerate}
    \item Since $n$ is odd, we can represent it in the form:$$n = 2k + 1$$where $k$ is any integer. Substitute $(2k + 1)$ for $n$ in the expression $n^2 - 1$:$$(2k + 1)^2 - 1 = 4k^2+4k = 4k(k+1)$$
    
    by law of excluded middle, \(k\) or \(k+1\) must be even. the product of any number with an even number must be even, thus \(k(k+1)\) must be even. Therefore \(k(k+1)\) can be written as \(2k_0 = k(k+1)\) for some \(k_0\). Together this means \(n^2-1 = 8k_0\) which means it is divisible by \(8\). 

    \item if \(n\) is odd, it can be written as \(2k+1\). 
    \begin{equation}
        (2k+1)^2+1 = 4k^2 + 4k + 2
    \end{equation}
    this number is clearly divisible by 2 and is thus even.
    \item 
    since \(n^2-1 = 8k_0\) and \(n^2+1 = 2k_1\) for some \(k_0,k_1\) in \(\mathbb Z\), the product can be written as \(16k_0k_1\) and thus \begin{equation}
        n^4-1 = (n^2 - 1)(n^2 + 1) = 16k_0k_1
    \end{equation}
    is divisible by 16.

    \item Let $n = 2k + 1$. We expand $(2k + 1)^4$ using the binomial theorem:$$(x + y)^n = \sum_{i=0}^{n} \binom{n}{i} x^i y^{n-i}$$
    For $n^4 = (2k + 1)^4$:\begin{equation}
        (2k)^4 + 4(2k)^3 + 6(2k)^2 + 4(2k)^1 + 1^4= 16k^4 + 32k^3 + 24k^2 + 8k + 1
    \end{equation}
    Now, we subtract 1 to get $n^4 - 1$:$$n^4 - 1 = 16k^4 + 32k^3 + 24k^2 + 8k$$

    \item We focus on the terms $a_2 k^2 + a_1 k$, which are $24k^2 + 8k$. We want to rewrite this in the form $bk^2 + c(k^2 + k)$. Notice that $24k^2$ can be split into $16k^2 + 8k^2$. \begin{equation}
        \begin{split}
            24k^2 + 8k &= 16k^2 + 8k^2 + 8k\\
            &= 16k^2 + 8(k^2 + k)
        \end{split}
    \end{equation} Here, $b = 16$ and $c = 8$.
    \item Since both terms are divisible by 8, the whole term is divisible by 8.
    \item \begin{equation}
    \begin{split}
        16k^4 + 32k^3 + 24k^2 + 8k &= 16k^4 + 32k^3 +16k^2 + 8k^2 + 8k\\
        &= 16(k^4 + 2k^3 +k^2) + 8(k^2 + k)\\
        &= 16(k^4 + 2k^3 +k^2) + 8(k+1)(k)\\
    \end{split}    \end{equation}
    since \((k+1)(k)\) is even, both terms are divisible by 16 and thus the original expression \(n^4-1\) is divisible by 16. 
\end{enumerate}



\item Prove or disprove:  $\forall x \in \mathbb{R}, \sqrt{x} \notin \mathbb{Q} \rightarrow x \notin \mathbb{Q} $.   \textbf{[5 points]}

\textbf{Sol}

This statement is false. A counterexample is \(\sqrt{2}\) which is not in \(\mathbb Q\), but \(2\) is. 

\item Prove or disprove: $ \sqrt[3]{3} \in \mathbb{Q}$.  \textbf{[5 points]} % $\forall x \in \mathbb{R}, x \notin \mathbb{Q} \rightarrow \sqrt[3]{x} \notin \mathbb{Q} $.  \textbf{[5 points]}

\textbf{Sol}

This statement is False.

Assume, for the sake of contradiction, that $\sqrt[3]{3}$ is a rational number. By definition, there exist two integers $p$ and $q$ ($q \neq 0$) such that:$$\sqrt[3]{3} = \frac{p}{q}$$
Assume that \(\frac{p}{q}\) is in its simplest form, i.e. \(\gcd(p,q) = 1\).

\begin{equation}
    \begin{split}
        \left(\sqrt[3]{3}\right)^3 &= \left(\frac{p}{q}\right)^3\\
        3 &= \frac{p^3}{q^3}
    \end{split}
\end{equation}

from this, we see that \(p^3\) is divisible by \(3\). Since \(3\) is prime, \(p\) must be divisible by 3 by Euler's Theorem. We can rewrite \(p\) as \(3p_1\). 
\begin{equation}
    (3p_1)^3 = 3q^3 = 27p_1^3
\end{equation}

which implies 
\begin{equation}
    q^3 = 9p_1^3 = 3(3p_1^3)
\end{equation}

This also means that \(q^3\) is divisible by \(3\), which applying Euler's Theorem we see implies \(q\) is also divisible by 3. 

\(q\) and \(p\) are divisible by 3 which is a contradiction since we assumed \(\gcd(q,p)=1\).



\item Prove:  If $n, m \in \mathbb{Z}$ and $7 n^4 + 11 m^3 = 214$ then $n$ and $m$ have the same parity.  \textbf{[5 points]}

\textbf{Sol}

Let us take the mod 2 of both sides

\begin{equation}
\begin{split}
        (7 n^4 + 11 m^3) \mod 2 &= (214 ) \mod 2 \\
        (n^4 + m^3) \mod 2 &= (0) \mod 2 \\
\end{split}\end{equation}

One can check exhausting all combinations of parity for \(n,m\), but a proof by contradiction is simpler.

Assume for the sake of contradiction that \(n,m\) have opposite parity. By induction, taking any number to some positive power retains its parity. Adding two numbers with opposite parity always results in parity 1 which is a contradiction.


\item Prove or disprove: If $x \notin \mathbb{Q}$, $y \in \mathbb{R}$, $x \neq 0$, and $y \neq 0$, then for all $n, m \in \mathbb{N}$,  $x^ny^m
%{\frac{1}{m}}  
\notin \mathbb{N}$.  \textbf{[5 points]}

\textbf{Sol}

This statement is Incorrect. A counterexample lies in \(x=\sqrt2, y=1, n=2, m=1\). (note, to prove that the forall is incorrect, we only need to find a single instance that fails the claim). 

\begin{equation}
    \sqrt2 ^2 \cdot 1^1 = 2 \in \mathbb N
\end{equation}


\item Consider the sets 
\[ S_1 = \{k^2+5: k \mbox{ is an even natural number} \},   S_2 = \{4k+5: k \in \mathbb{N} \}.    \]
\begin{enumerate}
	\item Prove $S_1 \subseteq S_2$. \textbf{[3 points]}
	\item Prove $\neg (S_2 \subseteq S_1)$. \textbf{[2 points]}
\end{enumerate}

\textbf{Sol}

\begin{enumerate}
    \item Let \(s_1\) be an arbitrary element in \(S_1\). Since \(k\) is even, it can be written as \(2k_0\) for some \(k_0\in\mathbb Z\). 
    
    \begin{equation}
        k^2+5 = (2k_0)^2+5 = 4k_0^2+5
    \end{equation}
    For any \(k_0\) in \(\mathbb Z\), the square of that number is a nonnegative integer, i.e. in \(mathbb N\). The form of this number matches with the specification in \(S_2\), thus it is an element in \(S_2\). Thus all elements in \(S_1\) are elements in \(S_2\),  i.e. \(S_1\subseteq S_2\). 

    \item To prove that \(S_2\) is not a subset of \(S_1\), we only need to find a single element in \(S_2\) that is not in \(S_1\). Consider \begin{equation}
        4(2)+5\in S_2
    \end{equation}

    this element is not in \(S_1\) since no even natural number squares to 8. (\(\sqrt{8} = 2\sqrt2\) is irrational.)
\end{enumerate}


\end{enumerate}

\end{document} 
